\input{preamble/preamble.tex}

\geometry{margin=0.85in}
\AtBeginDocument{\small}

\newcommand{\ExamNameField}{\noindent\textbf{Name:}\ \rule{0.7\linewidth}{0.4pt}\par\medskip}


\begin{document}
	\ExamTitleBlock{10th grade}{Learning evidence T3 Continuous probability and finance supplementary material solutions}{}
	
	\section*{Contents}
	\noindent\textbf{C1 Describes probabilities in continuous variables as a density function.}
	\begin{itemize}
		\item \hyperlink{c1-ex1}{Problem 1 — Rainfall Depth Model}
		\item \hyperlink{c1-ex2}{Problem 2 — Cooling Time Distribution}
	\end{itemize}
	\noindent\textbf{C2 Explains how a uniform random variable works.}
	\begin{itemize}
		\item \hyperlink{c2-ex1}{Problem 1 — Bus Arrival Wait Time}
		\item \hyperlink{c2-ex2}{Problem 2 — Random Placement in a Pool}
	\end{itemize}
	\noindent\textbf{C3 Expresses how to work with distributions that have a linear behaviour.}
	\begin{itemize}
		\item \hyperlink{c3-ex1}{Problem 1 — Battery Life with Increasing Density}
		\item \hyperlink{c3-ex2}{Problem 2 — Package Weights with Decreasing Density}
	\end{itemize}
	\noindent\textbf{C4 Interprets the concept of variance and standard deviation of a probability function.}
	\begin{itemize}
		\item \hyperlink{c4-ex1}{Problem 1 — Uniform Ticket Refunds}
		\item \hyperlink{c4-ex2}{Problem 2 — Triangular Delivery Times}
	\end{itemize}
	\noindent\textbf{C5 Uses the normal standard distribution and its tables to calculate probabilities in context situations.}
	\begin{itemize}
		\item \hyperlink{c5-ex1}{Problem 1 — Bottled Water Fill Volumes}
		\item \hyperlink{c5-ex2}{Problem 2 — Daily Study Time}
	\end{itemize}
	\noindent\textbf{C6 Employs the standardisation of the normal variable in real-life problems.}
	\begin{itemize}
		\item \hyperlink{c6-ex1}{Problem 1 — Test Scores Standardisation}
		\item \hyperlink{c6-ex2}{Problem 2 — Fruit Mass Quality Control}
	\end{itemize}
	\noindent\textbf{C7 Solves problems through the concept of inverse normal distribution.}
	\begin{itemize}
		\item \hyperlink{c7-ex1}{Problem 1 — Top 5\% Cable Lengths}
		\item \hyperlink{c7-ex2}{Problem 2 — Bottom 10\% Commute Times}
	\end{itemize}
	\noindent\textbf{C8 Summarizes the elements and characteristics that an investment portfolio may have.}
	\begin{itemize}
		\item \hyperlink{c8-ex1}{Problem 1 — Building a Student Savings Portfolio}
		\item \hyperlink{c8-ex2}{Problem 2 — Comparing Two Club Funds}
	\end{itemize}
	\noindent\textbf{C9 Establishes that the risk of a portfolio is determined by its probability distribution.}
	\begin{itemize}
		\item \hyperlink{c9-ex1}{Problem 1 — Two Portfolios with the Same Expected Return}
		\item \hyperlink{c9-ex2}{Problem 2 — Diversified vs. Concentrated Returns}
	\end{itemize}
	\noindent\textbf{C10 Evaluates finance-context situations with normal distribution and other continuous random variables.}
	\begin{itemize}
		\item \hyperlink{c10-ex1}{Problem 1 — Monthly Return Probabilities}
		\item \hyperlink{c10-ex2}{Problem 2 — Fees and Returns Together}
	\end{itemize}
	
	\ExamSection{C1 Describes probabilities in continuous variables as a density function.}
	
	\begin{ExamProblems}
		
		\hypertarget{c1-ex1}{}
		\item
		\subsection*{Problem 1 — Rainfall Depth Model}
		
		\textbf{Problem.}
		A town models the daily rainfall depth (in millimeters) on a rainy day with a continuous random variable
		\(X\) supported on \([0, 10]\). The density is defined as
		\[
		f(x)=k x,\quad 0\le x\le 10.
		\]
		
		\textbf{Question.} Find the constant \(k\), then compute the probability that the rainfall is between
		\(3\) mm and \(7\) mm. Interpret the result.
		
		\textbf{Solution.}
		A density must integrate to 1 over its support, so
		\[
		\int_0^{10} kx\,dx = 1.
		\]
		Compute the integral:
		\[
		k\left[\frac{x^2}{2}\right]_0^{10}=k\cdot\frac{100}{2}=50k.
		\]
		Thus \(50k=1\Rightarrow k=\frac{1}{50}.\)
		
		Now compute the requested probability:
		\[
		P(3\le X\le 7)=\int_3^7 \frac{1}{50}x\,dx
		=\frac{1}{50}\left[\frac{x^2}{2}\right]_3^7
		=\frac{1}{100}(49-9)=\frac{40}{100}=0.40.
		\]
		\textbf{Interpretation.} There is a \(40\%\) chance that rainfall on a rainy day is between 3 mm and 7 mm.
		
		% --------------------------------------------------
		
		\hypertarget{c1-ex2}{}
		\item
		\subsection*{Problem 2 — Cooling Time Distribution}
		
		\textbf{Problem.}
		A science class observes the cooling time (in minutes) of a heated metal rod. The random variable
		\(T\) is continuous on \([2, 8]\) with density
		\[
		f(t)=\frac{1}{12}(t-2),\quad 2\le t\le 8.
		\]
		
		\textbf{Question.} Verify that \(f(t)\) is a valid density, then find \(P(T>6)\). Interpret the result.
		
		\textbf{Solution.}
		Check the total area:
		\[
		\int_2^8 \frac{1}{12}(t-2)\,dt
		=\frac{1}{12}\left[\frac{(t-2)^2}{2}\right]_2^8
		=\frac{1}{12}\cdot\frac{36}{2}=\frac{1}{12}\cdot 18=1.
		\]
		So the density is valid.
		
		Now compute the probability:
		\[
		P(T>6)=\int_6^8 \frac{1}{12}(t-2)\,dt
		=\frac{1}{12}\left[\frac{(t-2)^2}{2}\right]_6^8
		=\frac{1}{12}\cdot\frac{36-16}{2}
		=\frac{1}{12}\cdot 10=\frac{5}{6}.
		\]
		\textbf{Interpretation.} About \(83\%\) of the time, cooling takes longer than 6 minutes.
		
	\end{ExamProblems}
	
	\ExamSection{C2 Explains how a uniform random variable works.}
	
	\begin{ExamProblems}
		
		\hypertarget{c2-ex1}{}
		\item
		\subsection*{Problem 1 — Bus Arrival Wait Time}
		
		\textbf{Problem.}
		A bus arrives uniformly at any time between 0 and 12 minutes after a student reaches the stop.
		Let \(W\) be the waiting time (minutes). Then \(W\sim\text{Uniform}(0,12)\).
		
		\textbf{Question.} Find \(P(W<5)\) and \(P(4\le W\le 9)\). Interpret both results.
		
		\textbf{Solution.}
		For a uniform distribution on \([a,b]\), the density is
		\[
		f(w)=\frac{1}{b-a}.
		\]
		Here \(a=0\), \(b=12\), so \(f(w)=\frac{1}{12}\).
		
		First probability:
		\[
		P(W<5)=\int_0^5 \frac{1}{12}\,dw=\frac{5}{12}.
		\]
		Second probability:
		\[
		P(4\le W\le 9)=\int_4^9 \frac{1}{12}\,dw=\frac{5}{12}.
		\]
		\textbf{Interpretation.} There is a \(\frac{5}{12}\) (about 41.7\%) chance the wait is under 5 minutes,
		and the same chance that the wait is between 4 and 9 minutes.
		
		% --------------------------------------------------
		
		\hypertarget{c2-ex2}{}
		\item
		\subsection*{Problem 2 — Random Placement in a Pool}
		
		\textbf{Problem.}
		A swimmer chooses a random depth along a sloped pool between 1 m and 3 m.
		Let \(D\sim\text{Uniform}(1,3)\).
		
		\textbf{Question.} Find \(P(D>2.5)\) and the expected depth \(E[D]\). Interpret both results.
		
		\textbf{Solution.}
		The density is \(f(d)=\frac{1}{3-1}=\frac{1}{2}\) for \(1\le d\le 3\).
		
		Probability:
		\[
		P(D>2.5)=\int_{2.5}^3 \frac{1}{2}\,dd=\frac{0.5}{2}=0.25.
		\]
		Expected value for a uniform \([a,b]\) is \(E[D]=\frac{a+b}{2}\), so
		\[
		E[D]=\frac{1+3}{2}=2.
		\]
		\textbf{Interpretation.} There is a 25\% chance the depth exceeds 2.5 m, and the average
		chosen depth is 2 m.
		
	\end{ExamProblems}
	
	\ExamSection{C3 Expresses how to work with distributions that have a linear behaviour.}
	
	\begin{ExamProblems}
		
		\hypertarget{c3-ex1}{}
		\item
		\subsection*{Problem 1 — Battery Life with Increasing Density}
		
		\textbf{Problem.}
		A battery's lifetime (hours) is modeled by a continuous random variable \(X\) on \([0, 6]\) with
		linear density
		\[
		f(x)=a+bx,\quad 0\le x\le 6,
		\]
		and \(f(0)=0\). That means the density starts at 0 and increases linearly.
		
		\textbf{Question.} Find \(f(x)\) explicitly, then compute \(P(X\le 3)\). Interpret the result.
		
		\textbf{Solution.}
		The condition \(f(0)=0\) gives \(a=0\), so \(f(x)=bx\).
		Total area must be 1:
		\[
		\int_0^6 bx\,dx=b\left[\frac{x^2}{2}\right]_0^6=b\cdot\frac{36}{2}=18b=1.
		\]
		Thus \(b=\frac{1}{18}\) and \(f(x)=\frac{1}{18}x\).
		
		Now compute
		\[
		P(X\le 3)=\int_0^3 \frac{1}{18}x\,dx=\frac{1}{18}\left[\frac{x^2}{2}\right]_0^3=\frac{1}{18}\cdot\frac{9}{2}=\frac{1}{4}.
		\]
		\textbf{Interpretation.} Only 25\% of batteries last 3 hours or less, so longer lives are more likely.
		
		% --------------------------------------------------
		
		\hypertarget{c3-ex2}{}
		\item
		\subsection*{Problem 2 — Package Weights with Decreasing Density}
		
		\textbf{Problem.}
		A delivery service models package weight (kg) by a linear density on \([0, 4]\):
		\[
		f(w)=a-bw,\quad 0\le w\le 4,
		\]
		with \(f(4)=0\) so heavier packages are less common.
		
		\textbf{Question.} Determine \(f(w)\), then find \(P(1\le W\le 3)\). Interpret the result.
		
		\textbf{Solution.}
		From \(f(4)=0\), we have \(a-4b=0\Rightarrow a=4b\).
		Total area is 1:
		\[
		\int_0^4 (a-bw)\,dw=1.
		\]
		Substitute \(a=4b\):
		\[
		\int_0^4 (4b-bw)\,dw=b\int_0^4 (4-w)\,dw
		=b\left[4w-\frac{w^2}{2}\right]_0^4
		=b(16-8)=8b.
		\]
		So \(8b=1\Rightarrow b=\frac{1}{8}\), and \(a=\frac{1}{2}\).
		Thus \(f(w)=\frac{1}{2}-\frac{1}{8}w\).
		
		Now compute
		\[
		P(1\le W\le 3)=\int_1^3 \left(\frac{1}{2}-\frac{1}{8}w\right)dw
		=\left[\frac{w}{2}-\frac{w^2}{16}\right]_1^3
		=\left(\frac{3}{2}-\frac{9}{16}\right)-\left(\frac{1}{2}-\frac{1}{16}\right)
		=\frac{15}{16}-\frac{7}{16}=\frac{1}{2}.
		\]
		\textbf{Interpretation.} About 50\% of packages weigh between 1 kg and 3 kg.
		
	\end{ExamProblems}
	
	\ExamSection{C4 Interprets the concept of variance and standard deviation of a probability function.}
	
	\begin{ExamProblems}
		
		\hypertarget{c4-ex1}{}
		\item
		\subsection*{Problem 1 — Uniform Ticket Refunds}
		
		\textbf{Problem.}
		A theater issues random refunds uniformly between \(\$0\) and \(\$20\) when a show is delayed.
		Let \(R\sim\text{Uniform}(0,20)\) denote the refund amount.
		
		\textbf{Question.} Compute \(E[R]\), \(\mathrm{Var}(R)\), and the standard deviation. Interpret the spread.
		
		\textbf{Solution.}
		For a uniform \([a,b]\):
		\[
		E[R]=\frac{a+b}{2},\quad \mathrm{Var}(R)=\frac{(b-a)^2}{12}.
		\]
		Here \(a=0\), \(b=20\), so
		\[
		E[R]=\frac{0+20}{2}=10,\quad \mathrm{Var}(R)=\frac{20^2}{12}=\frac{400}{12}=\frac{100}{3}.
		\]
		Standard deviation:
		\[
		\sigma=\sqrt{\frac{100}{3}}\approx 5.77.
		\]
		\textbf{Interpretation.} The average refund is \$10, and the typical variation is about \$5.77
		over or under the mean.
		
		% --------------------------------------------------
		
		\hypertarget{c4-ex2}{}
		\item
		\subsection*{Problem 2 — Triangular Delivery Times}
		
		\textbf{Problem.}
		Delivery time (hours) for a local service is modeled by
		\[
		f(t)=\frac{1}{8}t,\quad 0\le t\le 4.
		\]
		
		\textbf{Question.} Compute \(E[T]\), \(\mathrm{Var}(T)\), and the standard deviation. Interpret the result.
		
		\textbf{Solution.}
		First compute the mean:
		\[
		E[T]=\int_0^4 t\cdot \frac{1}{8}t\,dt=\frac{1}{8}\int_0^4 t^2\,dt
		=\frac{1}{8}\left[\frac{t^3}{3}\right]_0^4=\frac{1}{8}\cdot\frac{64}{3}=\frac{8}{3}.
		\]
		Next compute \(E[T^2]\):
		\[
		E[T^2]=\int_0^4 t^2\cdot \frac{1}{8}t\,dt=\frac{1}{8}\int_0^4 t^3\,dt
		=\frac{1}{8}\left[\frac{t^4}{4}\right]_0^4=\frac{1}{8}\cdot\frac{256}{4}=8.
		\]
		Variance is
		\[
		\mathrm{Var}(T)=E[T^2]-(E[T])^2=8-\left(\frac{8}{3}\right)^2
		=8-\frac{64}{9}=\frac{8}{9}.
		\]
		Standard deviation:
		\[
		\sigma=\sqrt{\frac{8}{9}}=\frac{2\sqrt{2}}{3}\approx 0.94.
		\]
		\textbf{Interpretation.} Average delivery time is \(\frac{8}{3}\) hours (about 2.67 h) with
		a typical spread of about 0.94 hours.
		
	\end{ExamProblems}
	
	\ExamSection{C5 Uses the normal standard distribution and its tables to calculate probabilities in context situations.}
	
	\begin{ExamProblems}
		
		\hypertarget{c5-ex1}{}
		\item
		\subsection*{Problem 1 — Bottled Water Fill Volumes}
		
		\textbf{Problem.}
		A bottling machine fills bottles with a mean of 500 mL and standard deviation of 8 mL.
		Assume fill volume \(X\sim N(500,8^2)\).
		
		\textbf{Question.} Use the standard normal table to find \(P(X<512)\). Interpret the result.
		
		\textbf{Solution.}
		Standardize:
		\[
		Z=\frac{X-\mu}{\sigma}=\frac{X-500}{8}.
		\]
		So
		\[
		P(X<512)=P\left(Z<\frac{512-500}{8}\right)=P(Z<1.5).
		\]
		From the standard normal table, \(\Phi(1.50)\approx 0.9332\).
		\textbf{Interpretation.} About 93.3\% of bottles are filled with less than 512 mL.
		
		% --------------------------------------------------
		
		\hypertarget{c5-ex2}{}
		\item
		\subsection*{Problem 2 — Daily Study Time}
		
		\textbf{Problem.}
		Daily study time for a group of students is approximately normal with mean 2.5 hours
		and standard deviation 0.6 hours. Let \(X\sim N(2.5,0.6^2)\).
		
		\textbf{Question.} Use the standard normal table to find \(P(2.0\le X\le 3.2)\).
		
		\textbf{Solution.}
		Standardize the endpoints:
		\[
		Z_1=\frac{2.0-2.5}{0.6}=-0.83\quad\text{(approx)},\qquad
		Z_2=\frac{3.2-2.5}{0.6}=1.17\quad\text{(approx)}.
		\]
		Thus
		\[
		P(2.0\le X\le 3.2)=P(-0.83\le Z\le 1.17)
		=\Phi(1.17)-\Phi(-0.83).
		\]
		From the table, \(\Phi(1.17)\approx 0.8790\) and \(\Phi(-0.83)=1-\Phi(0.83)\approx 1-0.7967=0.2033\).
		Therefore
		\[
		P(2.0\le X\le 3.2)\approx 0.8790-0.2033=0.6757.
		\]
		\textbf{Interpretation.} About 67.6\% of students study between 2.0 and 3.2 hours per day.
		
	\end{ExamProblems}
	
	\ExamSection{C6 Employs the standardisation of the normal variable in real-life problems.}
	
	\begin{ExamProblems}
		
		\hypertarget{c6-ex1}{}
		\item
		\subsection*{Problem 1 — Test Scores Standardisation}
		
		\textbf{Problem.}
		A math test score \(X\) is normally distributed with mean 70 and standard deviation 8.
		
		\textbf{Question.} Find the probability that a randomly chosen student scores above 85.
		Show the standardisation step clearly.
		
		\textbf{Solution.}
		Standardize:
		\[
		Z=\frac{X-70}{8}.
		\]
		Then
		\[
		P(X>85)=P\left(Z>\frac{85-70}{8}\right)=P(Z>1.875).
		\]
		From the standard normal table, \(\Phi(1.88)\approx 0.9699\).
		So
		\[
		P(Z>1.875)\approx 1-0.9699=0.0301.
		\]
		\textbf{Interpretation.} Only about 3\% of students score above 85.
		
		% --------------------------------------------------
		
		\hypertarget{c6-ex2}{}
		\item
		\subsection*{Problem 2 — Fruit Mass Quality Control}
		
		\textbf{Problem.}
		The mass of apples from an orchard is normally distributed with mean 180 g
		and standard deviation 12 g. Apples are considered "large" if they weigh more than 200 g.
		
		\textbf{Question.} Find the probability that an apple is classified as large.
		
		\textbf{Solution.}
		Standardize:
		\[
		Z=\frac{X-180}{12}.
		\]
		Compute
		\[
		P(X>200)=P\left(Z>\frac{200-180}{12}\right)=P(Z>1.67\text{ (approx)}).
		\]
		From the table, \(\Phi(1.67)\approx 0.9525\).
		Thus
		\[
		P(X>200)\approx 1-0.9525=0.0475.
		\]
		\textbf{Interpretation.} About 4.8\% of apples are large.
		
	\end{ExamProblems}
	
	\ExamSection{C7 Solves problems through the concept of inverse normal distribution.}
	
	\begin{ExamProblems}
		
		\hypertarget{c7-ex1}{}
		\item
		\subsection*{Problem 1 — Top 5\% Cable Lengths}
		
		\textbf{Problem.}
		Cables are manufactured with lengths \(L\sim N(50\text{ cm}, 1.2^2\text{ cm}^2)\).
		
		\textbf{Question.} Find the length \(\ell\) such that only 5\% of cables are longer than \(\ell\).
		State the target probability before using the inverse normal table.
		
		\textbf{Solution.}
		"Top 5\%" means
		\[
		P(L>\ell)=0.05\quad\Rightarrow\quad P(L\le \ell)=0.95.
		\]
		Convert to a standard normal value:
		\[
		P\left(Z\le \frac{\ell-50}{1.2}\right)=0.95.
		\]
		From the inverse normal table, \(z_{0.95}\approx 1.645\).
		So
		\[
		\frac{\ell-50}{1.2}=1.645\Rightarrow \ell=50+1.2(1.645)\approx 51.97\text{ cm}.
		\]
		\textbf{Interpretation.} About 5\% of cables are longer than 52.0 cm.
		
		% --------------------------------------------------
		
		\hypertarget{c7-ex2}{}
		\item
		\subsection*{Problem 2 — Bottom 10\% Commute Times}
		
		\textbf{Problem.}
		Commute times to school are normally distributed with mean 25 minutes and
		standard deviation 6 minutes.
		
		\textbf{Question.} Find the time \(t\) such that 10\% of commutes are shorter than \(t\).
		
		\textbf{Solution.}
		"Bottom 10\%" means
		\[
		P(T\le t)=0.10.
		\]
		Standardize:
		\[
		P\left(Z\le \frac{t-25}{6}\right)=0.10.
		\]
		From the inverse normal table, \(z_{0.10}\approx -1.28\).
		So
		\[
		\frac{t-25}{6}=-1.28\Rightarrow t=25-1.28(6)=25-7.68=17.32\text{ minutes}.
		\]
		\textbf{Interpretation.} About 10\% of students arrive in 17.3 minutes or less.
		
	\end{ExamProblems}
	
	\ExamSection{C8 Summarizes the elements and characteristics that an investment portfolio may have.}
	
	\begin{ExamProblems}
		
		\hypertarget{c8-ex1}{}
		\item
		\subsection*{Problem 1 — Building a Student Savings Portfolio}
		
		\textbf{Problem.}
		A student club invests \(\$10{,}000\) in three assets: a savings account (40\%), a government bond (35\%),
		and a stock fund (25\%). Expected annual returns are 2\%, 4\%, and 8\%, respectively.
		
		\textbf{Question.} Compute the expected portfolio return and describe the portfolio's
		elements (assets, weights, and risk characteristics) in words.
		
		\textbf{Solution.}
		Expected portfolio return is the weighted average:
		\[
		E[R]=0.40(0.02)+0.35(0.04)+0.25(0.08).
		\]
		Compute:
		\[
		E[R]=0.008+0.014+0.020=0.042=4.2\%.
		\]
		\textbf{Interpretation.} The portfolio mixes low-risk savings with medium-risk bonds and higher-risk
		stocks. The weights show more money in safer assets, so the expected return is moderate (4.2\%).
		It has diversification, some liquidity (savings), and some growth potential (stocks).
		
		% --------------------------------------------------
		
		\hypertarget{c8-ex2}{}
		\item
		\subsection*{Problem 2 — Comparing Two Club Funds}
		
		\textbf{Problem.}
		Two school clubs invest in different portfolios:
		\begin{itemize}
			\item Portfolio A: 60\% bonds, 40\% stocks (expected returns 4\% and 9\%).
			\item Portfolio B: 30\% bonds, 70\% stocks (same expected returns).
		\end{itemize}
		
		\textbf{Question.} Compute expected returns and summarize how the asset mix affects risk and return.
		
		\textbf{Solution.}
		Portfolio A expected return:
		\[
		E[R_A]=0.60(0.04)+0.40(0.09)=0.024+0.036=0.060=6.0\%.
		\]
		Portfolio B expected return:
		\[
		E[R_B]=0.30(0.04)+0.70(0.09)=0.012+0.063=0.075=7.5\%.
		\]
		\textbf{Interpretation.} Portfolio B has higher expected return because it holds more stocks, but
		stocks are typically more volatile. Portfolio A is more conservative with a larger bond share,
		so it has lower expected return and lower risk.
		
	\end{ExamProblems}
	
	\ExamSection{C9 Establishes that the risk of a portfolio is determined by its probability distribution.}
	
	\begin{ExamProblems}
		
		\hypertarget{c9-ex1}{}
		\item
		\subsection*{Problem 1 — Two Portfolios with the Same Expected Return}
		
		\textbf{Problem.}
		Two portfolios have the following annual return distributions:
		\[
		\begin{array}{lccc}
		\toprule
		\text{Return} & -4\% & 6\% & 16\% \\
		\midrule
		\text{Portfolio A} & 0.25 & 0.50 & 0.25 \\
		\text{Portfolio B} & 0.10 & 0.80 & 0.10 \\
		\bottomrule
		\end{array}
		\]
		
		\textbf{Question.} Compute the expected return and variance for each portfolio and decide
		which one is riskier.
		
		\textbf{Solution.}
		Expected return for A:
		\[
		E[R_A]=0.25(-0.04)+0.50(0.06)+0.25(0.16)= -0.01+0.03+0.04=0.06.
		\]
		Expected return for B:
		\[
		E[R_B]=0.10(-0.04)+0.80(0.06)+0.10(0.16)= -0.004+0.048+0.016=0.06.
		\]
		Both have the same expected return (6\%). Now compute variance.
		
		For A:
		\[
		E[R_A^2]=0.25(0.04^2)+0.50(0.06^2)+0.25(0.16^2)
		=0.25(0.0016)+0.50(0.0036)+0.25(0.0256)
		=0.0004+0.0018+0.0064=0.0086.
		\]
		\[
		\mathrm{Var}(R_A)=0.0086-(0.06)^2=0.0086-0.0036=0.0050.
		\]
		For B:
		\[
		E[R_B^2]=0.10(0.04^2)+0.80(0.06^2)+0.10(0.16^2)
		=0.10(0.0016)+0.80(0.0036)+0.10(0.0256)
		=0.00016+0.00288+0.00256=0.0056.
		\]
		\[
		\mathrm{Var}(R_B)=0.0056-(0.06)^2=0.0056-0.0036=0.0020.
		\]
		\textbf{Interpretation.} Portfolio A has a larger variance, so its distribution is more spread out
		and it is riskier, even though both have the same expected return.
		
		% --------------------------------------------------
		
		\hypertarget{c9-ex2}{}
		\item
		\subsection*{Problem 2 — Diversified vs. Concentrated Returns}
		
		\textbf{Problem.}
		A diversified portfolio (D) and a concentrated portfolio (C) have the following distributions:
		\[
		\begin{array}{lccc}
		\toprule
		\text{Return} & -8\% & 5\% & 18\% \\
		\midrule
		\text{Portfolio D} & 0.20 & 0.60 & 0.20 \\
		\text{Portfolio C} & 0.35 & 0.30 & 0.35 \\
		\bottomrule
		\end{array}
		\]
		
		\textbf{Question.} Compute the variance for each portfolio and explain how the distribution
		relates to risk.
		
		\textbf{Solution.}
		Compute expected returns:
		\[
		E[R_D]=0.20(-0.08)+0.60(0.05)+0.20(0.18)=-0.016+0.030+0.036=0.050.
		\]
		\[
		E[R_C]=0.35(-0.08)+0.30(0.05)+0.35(0.18)=-0.028+0.015+0.063=0.050.
		\]
		Both have expected return 5\%.
		
		Compute variances.
		For D:
		\[
		E[R_D^2]=0.20(0.08^2)+0.60(0.05^2)+0.20(0.18^2)
		=0.20(0.0064)+0.60(0.0025)+0.20(0.0324)
		=0.00128+0.00150+0.00648=0.00926.
		\]
		\[
		\mathrm{Var}(R_D)=0.00926-(0.05)^2=0.00926-0.0025=0.00676.
		\]
		For C:
		\[
		E[R_C^2]=0.35(0.08^2)+0.30(0.05^2)+0.35(0.18^2)
		=0.35(0.0064)+0.30(0.0025)+0.35(0.0324)
		=0.00224+0.00075+0.01134=0.01433.
		\]
		\[
		\mathrm{Var}(R_C)=0.01433-(0.05)^2=0.01433-0.0025=0.01183.
		\]
		\textbf{Interpretation.} Portfolio C has the larger variance because its distribution puts more
		probability on extreme returns. A wider distribution means higher risk.
		
	\end{ExamProblems}
	
	\ExamSection{C10 Evaluates finance-context situations with normal distribution and other continuous random variables.}
	
	\begin{ExamProblems}
		
		\hypertarget{c10-ex1}{}
		\item
		\subsection*{Problem 1 — Monthly Return Probabilities}
		
		\textbf{Problem.}
		A mutual fund's monthly return \(R\) is approximately normal with mean 1.2\% and
		standard deviation 3\%. Assume \(R\sim N(0.012,0.03^2)\).
		
		\textbf{Question.} Find the probability that the fund has a loss in a month
		(\(R<0\)). Interpret the result.
		
		\textbf{Solution.}
		Standardize:
		\[
		Z=\frac{R-0.012}{0.03}.
		\]
		Then
		\[
		P(R<0)=P\left(Z<\frac{0-0.012}{0.03}\right)=P(Z<-0.40).
		\]
		From the table, \(\Phi(-0.40)=1-\Phi(0.40)\approx 1-0.6554=0.3446\).
		\textbf{Interpretation.} There is about a 34\% chance of a negative return in a month.
		
		% --------------------------------------------------
		
		\hypertarget{c10-ex2}{}
		\item
		\subsection*{Problem 2 — Fees and Returns Together}
		
		\textbf{Problem.}
		An investment platform charges a random monthly fee \(F\) (in dollars) that is uniformly
		distributed between \(\$2\) and \(\$8\). The monthly portfolio return \(R\) (in dollars)
		is modeled by a normal distribution with mean \(\$25\) and standard deviation \(\$10\).
		
		\textbf{Question.} Find the probability that the net gain \(G=R-F\) exceeds \(\$20\).
		Clearly define the random variables and interpret the result.
		
		\textbf{Solution.}
		Define \(F\sim\text{Uniform}(2,8)\) with density \(f_F(f)=\frac{1}{6}\) on \([2,8]\),
		and \(R\sim N(25,10^2)\).
		We want
		\[
		P(G>20)=P(R-F>20).
		\]
		Condition on a fee value \(F=f\) and integrate over the fee distribution:
		\[
		P(G>20)=\int_2^8 P(R>20+f)\cdot \frac{1}{6}\,df.
		\]
		Standardize inside the integral:
		\[
		P(R>20+f)=P\left(Z>\frac{20+f-25}{10}\right)=P\left(Z>\frac{f-5}{10}\right).
		\]
		Thus
		\[
		P(G>20)=\frac{1}{6}\int_2^8 \left[1-\Phi\left(\frac{f-5}{10}\right)\right]df.
		\]
		Approximate using midpoint values (for Grade 10 estimation). At \(f=5\), \(Z=0\), so
		\(1-\Phi(0)=0.5\). At \(f=2\), \(Z=-0.3\) giving \(1-\Phi(-0.3)=\Phi(0.3)\approx 0.6179\).
		At \(f=8\), \(Z=0.3\), giving \(1-\Phi(0.3)\approx 0.3821\).
		Averaging these values gives about 0.50, so
		\[
		P(G>20)\approx 0.50.
		\]
		\textbf{Interpretation.} There is about a 50\% chance that the net gain exceeds \$20
		when both return and fee variability are considered.
		
	\end{ExamProblems}
	
\end{document}
